\slajdid{02-s024}
\begin{frame}[label=02-s024]
\gusttekst
Međutim zahvaljujući dualnosti logičkih kola, De Morganovim obrascima, moguće je doći i do još nekih rezultata. Ajde da vidimo prvo preko električnih šema. Da u postojećoj realizaciji „dodamo“ na izlaze logičkih kola prvog nivoa invertore odnosno kola prvog nivoa da u realizaciji pretvorimo u NI logička kola a u drugoj realizaciji u NILI logička kola. Da ne bi promenili funkciju na ulaze logičkih kola u drugom nivou ćemo takođe dodati invertore, odnosno napraviti ILI logičko kolo sa aktivnim logičkim nulama na ulazu u prvoj realizaciji (dualnost: ILI sa aktivnim logičkim nulama na ulazu jednako I sa aktivnom logičkom nulom na izlazu tj NI), dok ćemo dodavanjem invertora na ulaze I kola u drugoj realizaciji napraviti I logičko kola sa aktivnim logičkim nulama na ulazu (dualnost: I sa aktivnim logičkim nulama na ulazu jednako ILI sa aktivnom logičkom nulom na izlazu tj NILI)

\centering
\input{slike/tikz/s024_ni_i_invertovani_ulazi.tex}
\(\longrightarrow\)
\begin{tikzpicture}[font=\fontsize{8}{9.5}\selectfont,x=7mm,y=4mm,baseline=(current bounding box.center),
 DEinvertor/.style={not gate US,draw,scale=.6,minimum width=5mm,minimum height=4mm},
 DEprvinivo/.style={nand gate US,draw,scale=.6,logic gate inputs=nnn,rotate=-90,minimum width=5mm,minimum height=5mm},
 DEdrug nivo/.style={nand gate US,draw,scale=.6,logic gate inputs=nnn,rotate=-90,minimum width=5mm,minimum height=5mm}]
\def\DEkrajvodova{6.2}
% Gornji vod para je prava vrednost; donji je komplementna.
\foreach \y/\low/\lab in {6.5/5.9/C,5.3/4.7/B,4.1/3.5/A}{
 \node[DEinvertor] (n\lab) at (.65,\y) {};
 \node[DEinvertor] (m\lab) at (2.0,\y) {};
 \draw[DEvod] (0,\y) node[left] {\(\lab\)}--(n\lab.input);
 \draw[DEvod] (n\lab.output)--(m\lab.input);
 \draw[DEvod] (m\lab.output)--(\DEkrajvodova,\y);
 \coordinate (t\lab) at ($(n\lab.output)!.5!(m\lab.input)$);
 \draw[DEvod] (t\lab)--(t\lab |- 0,\low)--(\DEkrajvodova,\low);
}

\foreach \i/\x in {1/3.1,2/4.4,3/5.7}{
 \node[DEprvinivo] (g\i) at (\x,2.1) {};
}
\foreach \i/\j/\y in {1/3/5.9,1/2/5.3,1/1/4.1,2/3/6.5,2/2/4.7,2/1/4.1,3/3/6.5,3/2/5.3,3/1/3.5}{
 \draw[DEvod] (g\i.input \j)--(g\i.input \j |- 0,\y);
}
\node[DEdrug nivo] (z) at (4.4,.15) {};
\draw[DEvod] (g1.output)--(g1.output |- 0,.95)-|(z.input 3);
\draw[DEvod] (g2.output)--(z.input 2);
\draw[DEvod] (g3.output)--(g3.output |- 0,.95)-|(z.input 1);
\draw[DEvod] (z.output)--++(0,-.35) node[below] {\(F\)};

\end{tikzpicture}

\end{frame}
